\section{Algorithme exact : Séparation et Évaluation (Branch and Bound)}

Pour résoudre de manière exacte le problème d'Open Shop à 3 machines $O_3||\sum T_j$, nous implémentons un algorithme de séparation et évaluation (Branch and Bound). Cet algorithme explore implicitement toutes les combinaisons d'ordonnancement possibles tout en élaguant les branches sous-optimales. L'arbre de recherche possède une profondeur maximale de $3n$ (3 machines $\times$ $n$ jobs = $3n$ opérations).

\subsection{Représentation d'un n\oe ud}

Chaque n\oe ud de l'arbre représente un ordonnancement partiel. On maintient les informations suivantes :
\begin{itemize}
    \item $\text{planifie}[j][k] \in \{0, 1\}$ : 1 si l'opération du job $j$ sur la machine $k$ est planifiée, 0 sinon.
    \item $\text{fin\_machine}[k]$ : date de disponibilité de la machine $k$ (fin de la dernière opération planifiée sur $k$).
    \item $\text{fin\_job}[j]$ : date de disponibilité du job $j$ (fin de la dernière opération planifiée pour $j$).
    \item $\text{start}[j][k]$ : date de début de l'opération $(j, k)$ si elle est planifiée.
    \item $\text{profondeur}$ : nombre d'opérations déjà planifiées ($0 \leq \text{profondeur} \leq 3n$).
\end{itemize}

Le n\oe ud racine correspond à un ordonnancement vide ($\text{profondeur} = 0$, toutes les disponibilités à 0).

\subsection{Solution initiale et borne supérieure}

Pour initialiser efficacement l'algorithme, on utilise la meilleure solution trouvée par le VNS comme solution initiale $S_0$. Sa tardiveté totale fournit la borne supérieure initiale :
\begin{equation}
    UB = \sum_{j=1}^{n} T_j(S_0) = 10
\end{equation}
Toute branche dont la borne inférieure atteint ou dépasse $UB$ est immédiatement élaguée.

\subsection{Stratégie de séparation}

À chaque n\oe ud, les branches enfants correspondent aux choix possibles pour la prochaine opération à ordonnancer. L'ensemble des candidats est constitué de toutes les opérations $O_{jk}$ (job $j$ sur machine $k$) telles que :
\begin{itemize}
    \item $\text{planifie}[j][k] = 0$ (l'opération n'a pas encore été planifiée).
\end{itemize}

Dans un problème d'Open Shop, l'ordre de passage des machines pour un job est libre, ce qui signifie que toute opération non planifiée est candidate. La date de début d'une opération candidate est :
\begin{equation}
    \text{start}[j][k] = \max(\text{fin\_machine}[k],\; \text{fin\_job}[j])
\end{equation}

\subsection{Borne inférieure (évaluation)}

La borne inférieure $LB$ d'un n\oe ud estime le minimum de $\sum T_j$ atteignable en descendant dans cette branche. Pour chaque job $j$, on calcule une borne inférieure $LB_j$ sur sa tardiveté :

\textbf{Job terminé} (toutes les opérations planifiées) :
\begin{equation}
    C_j = \max_{k} \left( \text{start}[j][k] + p_{jk} \right), \qquad LB_j = \max(0,\; C_j - d_j)
\end{equation}

\textbf{Job incomplet} (il reste des opérations sur l'ensemble $R_j = \{k : \text{planifie}[j][k] = 0\}$) :
\begin{equation}
    C_j^{lb} = \max\left(
        \text{fin\_job}[j] + \sum_{k \in R_j} p_{jk},\;
        \max_{k \in R_j} \left( \text{fin\_machine}[k] + p_{jk} \right)
    \right)
\end{equation}
\begin{equation}
    LB_j = \max(0,\; C_j^{lb} - d_j)
\end{equation}

Le premier terme traduit le fait que les opérations restantes d'un job sont séquentielles (un job ne peut être que sur une machine à la fois). Le second terme traduit la contrainte de disponibilité des machines.

La borne inférieure globale est alors :
\begin{equation}
    LB = \sum_{j=1}^{n} LB_j
\end{equation}

\textbf{Règle d'élagage :} si $LB \geq UB$, la branche est coupée car elle ne peut pas améliorer la meilleure solution connue.

\subsection{Exploration en profondeur}

L'arbre est parcouru en profondeur d'abord, ce qui présente deux avantages :
\begin{itemize}
    \item \textbf{Mise à jour rapide de $UB$} : en descendant rapidement vers les feuilles, on découvre tôt des solutions complètes qui resserrent la borne supérieure.
    \item \textbf{Complexité spatiale réduite} : la mémoire nécessaire est proportionnelle à la profondeur de l'arbre ($O(3n)$), ce qui évite la saturation de la mémoire.
\end{itemize}

\subsection{Pseudocode}

\begin{algorithm}[H]
\caption{Branch and Bound --- exploration récursive}
\begin{algorithmic}[1]
\REQUIRE N\oe ud courant $N$, borne supérieure $UB$, meilleure solution $S^*$
\IF{$\text{profondeur} = 3n$}
    \STATE Calculer $\sum T_j$ de la solution complète
    \IF{$\sum T_j < UB$}
        \STATE $UB \leftarrow \sum T_j$, \quad $S^* \leftarrow$ solution courante
    \ENDIF
    \RETURN
\ENDIF
\FORALL{opérations $(j,k)$ telles que $\text{planifie}[j][k] = 0$}
    \STATE $\text{start}[j][k] \leftarrow \max(\text{fin\_machine}[k],\; \text{fin\_job}[j])$
    \STATE Planifier $(j,k)$ : mettre à jour $\text{fin\_machine}[k]$, $\text{fin\_job}[j]$ et $\text{profondeur} \leftarrow \text{profondeur} + 1$
    \STATE Calculer $LB$ du n\oe ud fils
    \IF{$LB < UB$}
        \STATE Appel récursif sur le n\oe ud fils
    \ELSE
        \STATE Élaguer (n\oe ud non exploré)
    \ENDIF
    \STATE Retour arrière : restaurer l'état précédent du n\oe ud
\ENDFOR
\end{algorithmic}
\end{algorithm}

\subsection{Déroulement}

\textit{Légende des couleurs : \textcolor{blue}{Job 1}, \textcolor{red}{Job 2}, \textcolor{green!70!black}{Job 3}.}

On part de la solution VNS avec $UB = 10$ et un n\oe ud racine vide. Au fil de l'exploration, $UB$ se resserre : dès qu'une solution à $\sum T_j = 3$ est trouvée, tout n\oe ud de borne $LB \geq 3$ est élagué, ce qui réduit fortement l'espace de recherche.

\subsubsection{Solution optimale trouvée}

L'algorithme explore 379 n\oe uds et en élague 1090, puis confirme l'optimalité de la solution $\sum T_j = 3$ :

\begin{figure}[H]
\centering
\begin{ganttchart}[
    vgrid, hgrid,
    x unit=1.2cm, y unit chart=0.7cm,
    bar/.append style={draw=black, line width=1pt}
]{1}{8}
\gantttitle{Temps}{8} \\
\gantttitlelist{1,...,8}{1} \\
\ganttbar[bar/.style={fill=red!40}]{$M_1$}{1}{1}
\ganttbar[bar/.style={fill=blue!40}]{}{3}{5}
\ganttbar[bar/.style={fill=green!40}]{}{6}{7} \\
\ganttbar[bar/.style={fill=blue!40}]{$M_2$}{1}{2}
\ganttbar[bar/.style={fill=red!40}]{}{3}{6}
\ganttbar[bar/.style={fill=green!40}]{}{8}{8} \\
\ganttbar[bar/.style={fill=green!40}]{$M_3$}{1}{5}
\ganttbar[bar/.style={fill=blue!40}]{}{6}{6}
\ganttbar[bar/.style={fill=red!40}]{}{7}{8}
\end{ganttchart}
\caption{Solution optimale B\&B --- $C_1=6$ ($T_1=0$), $C_2=8$ ($T_2=0$), $C_3=8$ ($T_3=3$) ; $\sum T_j = 3$}
\end{figure}

\subsection{Résultat}

\textbf{Statistiques d'exploration :} 379 n\oe uds explorés, 1090 n\oe uds élagués. L'algorithme de séparation et évaluation confirme que $\sum T_j = 3$ est la solution \textbf{optimale} (le tableau comparatif complet de toutes les méthodes figure en fin de section sur l'algorithme génétique). Le Job 3 contribue inévitablement une tardiveté de 3 (son temps total de traitement est $\sum_k p_{3k} = 8 > d_3 = 5$), tandis que les Jobs 1 et 2 terminent exactement à leur échéance. L'élagage par la borne inférieure permet de ne parcourir que 26\% des n\oe uds générés, démontrant l'efficacité de la borne combinée (contraintes séquentielles job + disponibilités machines).
